\chapter{Problem Definition and Formulation}
\label{chap:problem}

%% ============================================================
%% PROBLEM 1 — Automated ICD-10-CM Code Prediction (from Old Report)
%% ============================================================
\section{Problem 1: Automated ICD-10-CM Code Prediction}
\label{sec:p1:icd_problem}

\subsection{Scope of the Problem}
\label{sec:p1:scope}
Accurate labeling of clinical notes with \emph{International Classification of Diseases (ICD)} codes is recognized as a critical task in contemporary healthcare documentation. Clinical narratives, however, often manifest as lengthy, highly unstructured text that contains abbreviations, domain-specific terminology, and nuanced clinical indicators. A further complication arises when only a small fraction of ICD codes—here termed ``rare codes''—occurs so infrequently that very few labeled instances are available for machine learning models. These codes frequently remain under-detected, which can diminish overall coding accuracy and potentially undermine large-scale health analytics. The objective of the present work is to automate the assignment of \emph{multiple ICD codes} to each clinical document, with a specific focus on \emph{improving predictive effectiveness for rare labels}.

Failure to correctly classify clinical narratives or to assign all relevant labels can have far-reaching consequences. At the institutional level, coding errors may result in imprecise billing processes and resource misallocation. At the individual level, a missed or erroneous diagnosis code may obscure a patient's condition or hamper effective disease management. On a broader scale, large-scale analytics that rely on aggregated administrative data may not capture the true incidence or prevalence of rare medical conditions, causing gaps in research and healthcare policy.

\subsection{Central Challenges}
\label{sec:p1:central_challenges}
Although recent advances in machine learning have significantly improved automated ICD coding, several formidable obstacles persist:

\begin{itemize}
    \item \textbf{Long-Tail Distribution:} The majority of available data is dominated by a few frequent codes, leaving rare diagnoses severely underrepresented.
    \item \textbf{High-Dimensional Label Space:} Each clinical document may include multiple ICD codes, and semantic interdependencies among these labels complicate multi-label classification.
    \item \textbf{Textual Complexity:} The unstructured nature of clinical notes, featuring inconsistent formatting and heavy use of abbreviations and domain-specific jargon, complicates text processing.
    \item \textbf{Knowledge Integration:} Existing ontologies contain valuable domain knowledge, but straightforward classification pipelines often overlook hierarchical or semantic relationships among medical codes.
    \item \textbf{Explainability Needs:} Models designed for clinical applications must be comprehensible to end-users, who require transparency and trust when evaluating automatic code assignments.
\end{itemize}

These challenges pose general difficulties for all ICD codes, but the detrimental effects are exacerbated for rare labels. Systems that perform well for frequently observed diagnoses often fail to capture low-frequency conditions, an omission that motivates the development of strategies tailored to address class imbalance and explainability gaps.

\subsection{Mathematical Formulation of the Task}
\label{sec:p1:mathformulation}

A concise mathematical formulation is presented below to clarify how a trained model should transform unstructured clinical text into accurate ICD predictions.

\subsubsection{Defining the Space of Clinical Narratives}
Let \(\mathcal{X}\) represent the collection of all possible clinical narratives. Each element \(x \in \mathcal{X}\) is a textual record describing a patient's clinical episode. In practice, \(x\) may include:

\begin{itemize}
    \item \emph{Admission notes} reflecting the circumstances motivating a patient's hospitalization.
    \item \emph{Discharge summaries} documenting final diagnoses, interventions, and follow-up steps.
    \item \emph{Progress notes} covering daily clinical updates for extended hospital stays.
\end{itemize}

Despite a variety of internal structures and formats, each \(x\) can be treated as a string of tokens (or words) for classification purposes. Thus, the set \(\mathcal{X}\) defines the \emph{input space} over which the model operates.

\subsubsection{Defining the Label Space of ICD Codes}
Let \(\mathcal{L} = \{\ell_1, \ell_2, \dots, \ell_{|\mathcal{L}|}\}\) be the full set of ICD codes. The notation \(|\mathcal{L}|\) denotes the number of distinct codes, which can be large in real-world applications:

\begin{itemize}
    \item Major ICD revisions (ICD-9, ICD-10, ICD-11) contain thousands of codes.
    \item Many codes represent specialized conditions (subtypes or sub-phenotypes).
\end{itemize}

Since multiple diagnoses are commonly documented for a single admission, the label space is multi-label in nature, allowing each clinical document to be assigned any subset of \(\mathcal{L}\).

\subsubsection{Labeled Dataset Construction}
Assume that a set of \(N\) labeled clinical narratives is available:
\[
D = \{ (x_i, Y_i)\}_{i=1}^N,
\]
where \(x_i \in \mathcal{X}\) and \(Y_i \subseteq \mathcal{L}\). Each pair \((x_i, Y_i)\) consists of a discharge summary (or similar record) and the corresponding ICD codes assigned by professional coders. This dataset \(D\) forms the basis for supervised learning.

\subsubsection{Function Definition}
The goal is to learn a function
\begin{equation}
f: \mathcal{X} \;\longrightarrow\; 2^{\mathcal{L}},
\label{eq:p1:fundefinition}
\end{equation}
where \(f(x) = \hat{Y}\) for \(\hat{Y} \subseteq \mathcal{L}\). Here, \(2^{\mathcal{L}}\) represents the power set of \(\mathcal{L}\). For any clinical document \(x\), the prediction \(\hat{Y}\) is the subset of codes that the model deems relevant. In principle, \(\hat{Y}\) can be empty, although that outcome is uncommon in practical medical scenarios.

In most current architectures, a multi-label classification approach is adopted. Specifically, the model generates a probability or score for each code in \(\mathcal{L}\), and thresholding (or a related method) is then applied to map those scores to a predicted subset.

\subsubsection{Objective: Composite Metric Maximization}
A composite performance metric \(\Omega\) is introduced to reflect the importance of both overall correctness and the accurate identification of rare labels. In formal terms, the aim is:

\begin{equation}
\underset{f}{\text{maximize}}\;\;\; \Omega(\text{Accuracy of } f,\;\text{Rare-Codes Metrics}).
\label{eq:p1:objomega}
\end{equation}

Depending on institutional priorities, \(\Omega\) may combine measures such as micro-F1 or micro-AUC (evaluating overall performance across all codes) with macro-F1 or a specialized metric that focuses on rare codes. A weighted sum or a multi-objective formulation can be employed. This dual emphasis on both frequent and rare diagnoses ensures that gains do not come at the expense of clinically significant but infrequent conditions.

\subsection{Significance of Rare Codes}
\label{sec:p1:rare_codes}

Rare codes pose a particular challenge due to limited representation in real-world data, yet these codes often signal clinically important conditions.

\subsubsection{Definition and Challenges}
A rare code may appear only a few times in a large dataset. Let \(\mathcal{L}_{\mathrm{rare}} \subseteq \mathcal{L}\) be the set of codes whose frequencies fall below a chosen threshold \(\tau\):
\[
\mathcal{L}_{\text{rare}} 
= 
\{\ell \in \mathcal{L} \,\mid\, \text{Frequency}(\ell) < \tau\}.
\]

Typical values for \(\tau\) might be a fixed count (e.g., 50) or a percentile (e.g., the bottom 10\%). Learning to detect rare codes is complicated by these labels' scarcity in training data and the specialized language often observed in notes that reference such conditions. Overall loss functions tend to favor frequent labels, marginalizing rare ones.

\subsubsection{Impact of Rare Code Misclassification}
Misclassification of rare diagnoses can produce substantial adverse effects:

\begin{itemize}
    \item \textbf{Clinical Impact:} A rare but critical diagnosis, if missed by an automated system, may affect patient care when clinicians rely on coded data for decision support.
    \item \textbf{Research Bias:} Underrepresentation of rare conditions in large-scale analyses can mask the true disease burden.
    \item \textbf{Billing and Compliance:} Inaccurate coding for rare conditions can lead to underpayment or claim denial, as well as potential auditing complexities.
\end{itemize}

Consequently, rare code detection forms a key component of robust automated coding systems.

\subsubsection{Recall and Precision Constraints for Rare Codes}
For a clinical note \(x\) with true codes \(Y\), where \(Y \cap \mathcal{L}_{\text{rare}} \neq \emptyset\), recall and precision for the rare portion are of particular concern. Low recall implies a failure to capture rare conditions, whereas low precision may yield an excess of false positives for these labels. Either error type can be problematic in a medical context, emphasizing the importance of balanced metrics.

\subsubsection{Analytical Example of Class Imbalance}
A simple numerical illustration can clarify the issue. Consider 5,000 labeled notes, and suppose:
\begin{itemize}
    \item A frequent code \(\ell_{\mathrm{common}}\) appears in 1,200 notes.
    \item A rare code \(\ell_{\mathrm{rare}}\) appears in 20 notes.
\end{itemize}
A classifier that never predicts \(\ell_{\mathrm{rare}}\) incurs only 20 false negatives on a base of 5,000, so the impact on the global accuracy may be less than 1\%. Such a result can inadvertently encourage ignoring rare codes. This outcome underscores the need for targeted strategies addressing low-frequency classes.

\subsection{Proposed Methodological Focus}
\label{sec:p1:methodfocus}
Several methodological approaches are planned to mitigate performance deficits in both frequent and rare ICD codes:

\begin{enumerate}
    \item \textbf{Data Augmentation for Rare Codes:} Synthetic data generation or specialized oversampling strategies that boost representation of low-frequency labels.
    \item \textbf{Advanced Multi-Label Architectures:} Network designs that incorporate hierarchical embeddings and structured label representations, taking advantage of semantic relationships within the ICD taxonomy.
    \item \textbf{Integration of External Knowledge:} Leveraging medical ontologies (e.g., SNOMED, UMLS) to enhance label embeddings and encourage clinically consistent classification.
    \item \textbf{Model Interpretability:} Developing explainable frameworks where predicted ICD codes can be traced back to their textual justifications, thereby fostering user trust.
\end{enumerate}

Each of these elements seeks to address the data imbalance inherent in ICD coding tasks while preserving or improving performance on all relevant labels, including those with scarce training instances.

\subsection{Evaluation Framework}
\label{sec:p1:evalframework}
Evaluation in this domain requires a range of measures to capture overall performance and rare-code outcomes:

\begin{itemize}
    \item \textbf{Micro-F1}: Reflects globally aggregated precision and recall. 
    \item \textbf{Macro-F1}: Averages precision and recall across codes, thereby placing additional emphasis on rare labels.
    \item \textbf{Recall@k}: Indicates how effectively the model recovers correct codes in its top-\(k\) predictions, a practical measure for coding assistance tools.
    \item \textbf{Rare-Code Metrics}: Tracks performance (e.g., precision, recall, F-score) on the subset \(\mathcal{L}_{\mathrm{rare}}\), ensuring improvements are realized for both frequent and rare labels.
\end{itemize}

A balanced presentation of both micro- and macro-level metrics is essential, given the skewed distribution of ICD codes. It is possible for a model to exhibit strong micro-F1 yet fail to detect many rare diagnoses, which would be revealed by a lower macro-F1 or poor rare-code metrics.

\bigskip
\noindent\textbf{Note on Problem~1 status.} The methodology for automated ICD-10-CM code prediction was proposed and the problem was formally defined as described above; however, full model training and synthetic data generation were \emph{not completed} due to high computational infrastructure costs. Preliminary results (baseline model evaluation and data exploration) are reported in Chapter~\ref{chap:results}. The complete execution of this pipeline remains a future work item (see Chapter~\ref{chap:conclusion}).


%% ============================================================
%% PROBLEM 2 — SA-AKI Early Mortality Risk Stratification (Current Report, preserved verbatim)
%% ============================================================
\section{Problem 2: SA-AKI Early Mortality Risk Stratification}
\label{sec:saaki_problem}

\noindent\textbf{Terminology note.} The SA-AKI research task is internally organized into three sub-problems: an unsupervised exploratory analysis (Sub-problem~A), a supervised classification task (Sub-problem~B), and a survival analysis task (Sub-problem~C). These correspond to \emph{sub-problems within Problem~2} and are \textbf{distinct} from the thesis-level Problem~1 (ICD code prediction) and Problem~2 (SA-AKI mortality) framing.

\bigskip

\subsection{Clinical decision and scope}
Sepsis–associated acute kidney injury (SA--AKI) constitutes a high–risk state coupling infection, shock physiology, and acute renal dysfunction. The \emph{decision point} is fixed: at the end of the first ICU day ($T_{\mathrm{obs}}{=}24$~h after ICU admission), the risk of in–hospital death is to be estimated for patients meeting SA--AKI criteria using only information available by $T_{\mathrm{obs}}$. The primary goal was to develop a predictive model that is both \emph{discriminative} and \emph{calibrated}. Two auxiliary objectives support this decision: (i) systematic \emph{exploratory data analysis} to discover clinical phenotypes in the early feature space, and (ii) \emph{time–to–event} modeling that translates the same early covariates into survival functions for \emph{time to death}. All cohort rules adhere to Sepsis--3 for sepsis and KDIGO for AKI, and reporting follows the TRIPOD statement.

\noindent\textbf{Out of scope.} This study does not estimate causal treatment effects or recommend specific therapies. The problem is predictive; models summarize risk conditional on the information set $\mathcal{I}(T_{\mathrm{obs}})$.

\subsection{Observational unit, clocks, and information sets}
Each observational unit corresponds to the \emph{first} eligible ICU admission for a patient in the SA--AKI cohort. Let ICU admission be $t{=}0$ and define the early feature window as
\[
[0,\,T_{\mathrm{obs}}], \qquad T_{\mathrm{obs}}=24~\text{hours}.
\]
All time–varying signals in this window—vital signs, laboratory measurements, fluids/urine output, and early therapies—are summarized by deterministic, leakage–safe functionals into a feature vector
\[
x_i \in \mathbb{R}^p,
\]
augmented with static demographics and comorbidity indicators. For the binary classification task, a one–time stratified 80/20 train–test split yielded training and test sets with $(8028, 162)$ and $(2008, 162)$ dimensions, respectively (including the target column). The positive (death) rate in the training set was $26.74\%$. The survival analysis used the same patient split with a feature matrix of $(8028, 124)$ and $(2008, 124)$ before the inclusion of time and event columns.\footnote{These dataset sizes and event fractions reflect the final data preparation used in model building and survival analysis, as detailed in the notebooks.}
The available information for patient $i$ at the decision time is the $\sigma$–field
\[
\mathcal{I}_i(T_{\mathrm{obs}})=\{\text{EHR observations with time stamp}\le T_{\mathrm{obs}}\}.
\]
\emph{Temporal discipline.} Any statistic requiring $t>T_{\mathrm{obs}}$ (e.g., total ICU length of stay, whole–stay extrema/means, post–baseline therapies) is deemed inadmissible for prediction and is excluded by a scripted leakage audit.

\paragraph{Target population and case definition.}
Sepsis is operationalized by Sepsis--3 (suspected infection plus a daily SOFA score increase of $\ge 2$); AKI is defined by KDIGO serum–creatinine and urine–output criteria. SA--AKI requires AKI onset within $\pm48$~h of sepsis onset. Only adults ($\ge$18~y) with an ICU length of stay $\ge$24~h, and only the first ICU stay per patient, are included. Cohort definition timestamps $t_{\text{sepsis}}$ and $t_{\text{aki}}$ serve only to identify the cohort and \emph{are not used as predictors}.

\subsection{Data–generating view and notation}
Let $(X,Y,T^{\mathrm{death}},\delta^{\mathrm{death}})\sim P$ on
\[
\mathcal{X}\times\{0,1\}\times\mathbb{R}_+\times\{0,1\},
\]
where $X$ is the T0--T24 feature vector; $Y$ is in–hospital death (1/0); $T^{\mathrm{death}}$ is time from ICU admission to in–hospital death or censoring at discharge; and $\delta^{\mathrm{death}}$ is the event indicator (1=death, 0=censored). Independent, identically distributed samples $\{(x_i,\cdot)\}_{i=1}^n$ are observed. A one–time \emph{80/20} split yields $(\mathcal{D}_{\mathrm{train}},\mathcal{D}_{\mathrm{test}})$. Any internal validation split for tasks like hyperparameter tuning or calibration is carved out exclusively \emph{from within} $\mathcal{D}_{\mathrm{train}}$. The survival analysis reuses the same patient split with right censoring.

\subsection{Sub-problem A (Unsupervised): Exploratory Analysis and Phenotyping}
\textbf{Goal.} To systematically explore the early feature space of the training data to identify clinically relevant patient phenotypes without using outcome labels. This involves analyzing the distributions of key biomarkers, assessing correlations between features and the mortality outcome, and comparing clinical characteristics (e.g., vital signs, lab values, SOFA components) between patients who survived and those who did not. This exploratory phase informs feature importance and provides clinical context for the supervised models.

\subsection{Sub-problem B (Supervised, Primary): Early–Window Mortality Risk}
\subsubsection{Label and objective}
Define the binary label
\[
Y=\mathbb{I}\{\text{in–hospital death prior to discharge}\}.
\]
A probabilistic classifier $f_{\theta}:\mathcal{X}\to[0,1]$ is sought such that
\[
\hat p(x)=f_{\theta}(x)\approx \Pr(Y=1\mid X=x,\ \mathcal{I}(T_{\mathrm{obs}})).
\]
The learning criterion on training data is the average negative log–likelihood (cross-entropy),
\[
\hat\theta\in\arg\min_{\theta}\;\frac{1}{m}\sum_{j=1}^{m}
\big\{-y_j\log f_{\theta}(x_j)-(1-y_j)\log(1-f_{\theta}(x_j))\big\},
\]
optimized using cross–validation within $\mathcal{D}_{\mathrm{train}}$.

\subsubsection{Performance targets}
Two properties are co–primary:
\begin{enumerate}[label=(\alph*)]
  \item \textbf{Discrimination:} High Area Under the Receiver Operating Characteristic Curve (AUROC) and F1-score on the held-out test set.
  \item \textbf{Calibration:} Reliable probabilities, assessed by generating a final calibrated model and evaluating its performance.
\end{enumerate}
Calibration is performed on a validation split drawn from $\mathcal{D}_{\mathrm{train}}$ using isotonic regression; final probabilities are $\tilde p=c\circ f_{\theta}(x)$, where $c$ denotes the calibrator.

\subsubsection{Decision rule and operating point}
Clinical actions are thresholded via $\delta_{\tau}(x)=\mathbb{I}\{\tilde p(x)\ge \tau\}$. The optimal operating point $\tau$ is selected on an internal validation split by maximizing Youden's $J$ index. The final model's utility across a range of thresholds is also assessed using decision-curve analysis.

\subsubsection{Hypothesis classes and regularization}
A spectrum of models, $\mathcal{H}$, is benchmarked to assess performance across different functional capacities: Logistic Regression (with and without class weights, and with various resampling techniques), regularized models (Ridge, Lasso, ElasticNet), Linear and Quadratic Discriminant Analysis (LDA/QDA), Gaussian Naive Bayes, Decision Tree, Random Forest, gradient-boosted trees (XGBoost, LightGBM), and a deep neural network (DNN++). Hyperparameters for tree-based models are tuned using the Optuna framework within training folds.

\subsubsection{Interpretability requirement}
For the final selected model ($f_{\theta}$), feature contributions are explained using SHapley Additive exPlanations (SHAP). This provides both global feature importance rankings and local, instance-level explanations for individual patient predictions, satisfying the need for model transparency.

\paragraph{Constraints and guardrails.}
(i) \emph{Temporal:} All predictive features are derived strictly from the $[0,24]$\,h window post-ICU admission. (ii) \emph{Independence:} Only the first ICU stay per patient is included. (iii) \emph{Siloed preprocessing:} All data transformations (imputation, scaling, indicator creation) are learned on the training set and applied without modification to the test set. (iv) \emph{No peeking:} The test set is held out and used only for a single final evaluation. (v) \emph{Cohort times not used as predictors:} $t_{\text{sepsis}}$ and $t_{\text{aki}}$ are used only for case definition.

\subsection{Sub-problem C (Survival): Early–Covariate Time to Death}
Using the same T0--T24 covariates, the \emph{time to in–hospital death} is modeled with right censoring at hospital discharge. For the event time data $(T^{\mathrm{death}},\delta^{\mathrm{death}})$, the conditional hazard function $h(t\mid x)$ is estimated via:
\begin{itemize}
  \item \textbf{Cox Proportional Hazards (PH):} A semi-parametric model specified as $h(t\mid x)=h_0(t)\exp\{\beta^\top x\}$, implemented with the \texttt{lifelines} library. Proportional hazards assumptions are checked.
  \item \textbf{DeepSurv:} A neural network-based Cox model where $h(t\mid x)=h_0(t)\exp\{g_{\psi}(x)\}$, with $g_\psi$ parameterized by a custom deep network implemented in PyTorch.
\end{itemize}
Performance is evaluated using Harrell's $C$–index, time-dependent AUC, calibration plots at clinically relevant horizons (e.g., 90, 180, 365 hours), and Kaplan-Meier curve stratification by predicted risk.

\subsection{Evaluation protocol}
For the classification task, a one–time 80/20 train–test split is used. An internal validation split is carved from the training set for model selection, calibration, and thresholding. The frozen pipeline is then applied once to the test set. For the survival task, the same patient partitions are used. Primary test metrics for the classifier include AUROC, F1-score, precision, and recall. For the survival models, primary metrics are Harrell's $C$-index, calibration, and risk stratification confirmed by the log-rank test.

\subsection{Assumptions and risks}
\begin{itemize}
  \item \textbf{Availability:} All predictors are assumed to be observable by $T_{\mathrm{obs}}$; features with excessive missingness or not typically available early are excluded.
  \item \textbf{Measurement:} Units are harmonized and physiological plausibility is enforced. The missing data mechanism is assumed to be Missing at Random (MAR) after finding it is not Missing Completely at Random (MCAR), justifying the use of imputation and missingness indicators.
  \item \textbf{Domain:} Models are developed on the MIMIC--IV dataset; transportability to external centers may require recalibration.
  \item \textbf{Predictive, not causal:} The models identify predictive associations, not causal relationships.
\end{itemize}

\subsection{Compact statement of aims}
\noindent\textbf{Aim\,A (Phenotyping).} To conduct a thorough exploratory data analysis on the T0--T24 training data to identify and visualize the key clinical characteristics and phenotypes that differ between survivors and non-survivors.

\noindent\textbf{Aim\,B (Primary classifier).} To train and validate a calibrated probabilistic model for in–hospital death, $\tilde p(x)=c\circ f_{\theta}(x)$, ensuring high discrimination and utility through a strict, leakage-controlled pipeline.

\noindent\textbf{Aim\,C (Survival).} Using the same early covariates, to estimate the survival function $S^{\mathrm{death}}(t\mid x)$ with Cox and DeepSurv models and evaluate their performance via C-index, calibration, and risk stratification.
